
\section{Method}\label{sec:method}

We argue that layout understanding and basic math operation capabilities are the key elements for performing visual language understanding/reasoning. We inject such capabilities to the model by proposing two pretraining tasks: \textbf{chart derendering} (\Cref{sec:plot_derendering}) and \textbf{math reasoning} (\Cref{sec:math_reasoning}) which we introduce in detail in the following sections.

\subsection{Chart Derendering}\label{sec:plot_derendering}
Plots and charts are usually generated by an underlying data table and a piece of code. Code decides the overall layout of the figure (e.g., type, direction, color/shape scheme of the chart) and the underlying data table decides the actual numbers and the groupings of them. Both the data and code are sent to a compiler/rendering engine to create the final image. 
To understand a chart one needs to discover the visual patterns in the image, effectively parse and group them to extract the key information. Reversing the plot rendering process demands all such capabilities and can thus serve as a perfect pretraining task.

In practice, it is challenging to simultaneously obtain charts, their underlying data tables, and their rendering code. To collect sufficient pretraining data, we independently accumulate (chart, code) and (chart, table) pairs. For (chart, code), we crawl all GitHub IPython notebooks with appropriate licenses and extract blocks with figures. A figure and the code block right before it are saved as a (chart, code) pair.\footnote{Note that the code snippet can be noisy since earlier blocks could also be relevant for generating the figure and also the snippet may contain bits of code that is irrelevant to generating the figure. Also note that the data table is frequently missing and usually not hardcoded in the notebook. As a result, we collect (chart, table) pairs separately.} For (chart, table) pairs, we explored two sources. First is to manually write code for converting web-crawled Wikipedia tables from ~\citet{herzig-etal-2020-tapas} to charts.
%around 200 lines of Python code, 
We randomly combine several plotting options. The key random variables include: using either \texttt{matplotlib} or \texttt{seaborn} as the plotting package; using either bar, line, or pie chart; styles and colors of the charts; whether to show numbers explicitly on the graph; font and size of the texts. Besides our own synthetic data, we also add chart-table pairs generated by %\citet{masry-etal-2022-chartqa} and
\citet{methani2020plotqa} (from PlotQA) to diversify the pretraining corpus. The second source is web-crawled chart-table pairs. Websites such as Statista provides both. We directly use the chart-table pairs crawled by \citet{masry-etal-2022-chartqa} (from ChartQA), containing around 20k pairs in total from four websites: Statista, Pew, Our World in Data, and OECD.\footnote{See \Cref{sec:dataset_details} for links.}

Note that to avoid leaking test information for the PlotQA and ChartQA tasks which use the same chart data as pretraining, we only use the chart-table pairs in the training sets for pretraining and test tables/charts are strictly excluded. In ablation study (\Cref{sec:ablations}), we will show that chart-table from both sources are useful and having a diverse set of chart-table pairs is always better. However, using only our synthetic data brings very significant improvement already, suggesting that the concept of chart derendering can be easily transferred to charts of other domains (including real-world charts).

%TODO: show some pretraining examples for the tasks


\subsection{Math Reasoning}\label{sec:math_reasoning}
Reasoning over visual language requires (1) effective recognition and grouping of the visual elements and also (2) applying mathematical operations (such as sorting, min/max, averaging, etc.) on top of them. Plot derendering addresses (1) but (2) is still lacking in the current pretraining framework. As a result, we propose to explicitly inject numerical reasoning knowledge to the image-to-text model by learning from textual math datasets.

We use two existing textual math reasoning datasets, MATH \citep{saxton2019analysing} and DROP \citep{dua-etal-2019-drop} for pretraining. MATH is synthetically created, containing two million training examples per module (type) of questions (see \Cref{sec:dataset_details} for a comprehensive listing of modules included in \model{} pretraining).
DROP is a reading-comprehension-style QA dataset where the input is a paragraph context and a question. DROP has 96k question and answer pairs over 6.7K paragraphs.\footnote{Note that for all datasets used for pretraining, we always use only the training set if there exists a split.} To solve questions in DROP, the model needs to read the paragraph, extract relevant numbers and perform numerical computation to predict the answer. We found both datasets to be complementarily helpful. MATH contains large amounts of questions and is categorized which helps us identify math operations needed to explicitly inject to the model. DROP's reading-comprehension format resembles the typical QA format where models need to simultaneously perform information extraction and reasoning. In practice, we render inputs of both datasets into images (concatenating the context and question for DROP). The image-to-text model is trained to decode the answer given the redered image. Examples of MATH and DROP can be found in \Cref{fig:front_page} (in light red).

Besides the two newly proposed pretraining strategies, to prevent catastrophic forgetting, we also keep applying the screenshot parsing pretraining from Pix2Struct \citep{lee2022pix2struct}. Specifically, given screenshot of a website (where parts of the website is masked), the image-to-text transformer needs to predict the underlying simplified HTML code that could render the original unmasked website screenshot. The final pretraining task is a mixture of all aforementioned tasks. We discuss the mixture weights in \Cref{sec:exp_setups}.
